Treating the ODE solver as a first-class experimental variable rather than a black box reveals
structure that a black-box treatment cannot: solver-induced bias in recovered SIR parameters obeys
the solver's own convergence order almost exactly, and this bias competes with noise-induced
variance in a way that can be located precisely as a crossover noise level $\sigma^*$. Below
$\sigma^*$, the numerical-analysis choices this course teaches -- which solver, what step size --
are the dominant factor in whether an epidemiological parameter estimate is trustworthy. Above it,
they are not, and the modeler's effort is better spent on more data or a better noise model.
Reaching this conclusion used the methods of all four CSE 402 labs and the optimization home
assignment (Table~\ref{tab:course}): root finding and LU factorization inside the gold standard
and the implicit solvers, the power method for every eigenvalue and condition number,
least-squares optimization for estimation, quadrature and
interpolation for the reference and observation operators, and random number generation with
Monte Carlo simulation for every noise experiment. Every
figure and number in this report is generated by a single reproducible pipeline
(\texttt{make all}) from the code in this repository.
